\documentclass[10pt,a4paper]{amsart}
\usepackage{amsmath,amssymb,amsthm,graphicx}
% For commutative diagrams we use Xy-pic
\usepackage[all]{xy}

% We use 2cell for 2-commutative diagrams.
\xyoption{2cell}
% Theorem environments.
%
\theoremstyle{plain}
\newtheorem{theorem}[subsection]{Theorem}
\newtheorem{proposition}[subsection]{Proposition}
\newtheorem{lemma}[subsection]{Lemma}

\theoremstyle{definition}
\newtheorem{definition}[subsection]{Definition}
\newtheorem{example}[subsection]{Example}
\newtheorem{exercise}[subsection]{Exercise}
\newtheorem{situation}[subsection]{Situation}

\theoremstyle{remark}
\newtheorem{remark}[subsection]{Remark}
\newtheorem{remarks}[subsection]{Remarks}

\numberwithin{equation}{subsection}

% Macros
%
\def\lim{\mathop{\mathrm{lim}}\nolimits}
\def\colim{\mathop{\mathrm{colim}}\nolimits}
\def\Spec{\mathop{\mathrm{Spec}}}
\def\Hom{\mathop{\mathrm{Hom}}\nolimits}
\def\Ext{\mathop{\mathrm{Ext}}\nolimits}
\def\SheafHom{\mathop{\mathcal{H}\!\mathit{om}}\nolimits}
\def\SheafExt{\mathop{\mathcal{E}\!\mathit{xt}}\nolimits}
\def\Sch{\mathit{Sch}}
\def\Mor{\mathop{\mathrm{Mor}}\nolimits}
\def\Ob{\mathop{\mathrm{Ob}}\nolimits}
\def\Sh{\mathop{\mathit{Sh}}\nolimits}
\def\NL{\mathop{N\!L}\nolimits}
\def\CH{\mathop{\mathrm{CH}}\nolimits}
\def\proetale{{pro\text{-}\acute{e}tale}}
\def\etale{{\acute{e}tale}}
\def\QCoh{\mathit{QCoh}}
\def\Ker{\mathop{\mathrm{Ker}}}
\def\Im{\mathop{\mathrm{Im}}}
\def\Coker{\mathop{\mathrm{Coker}}}
\def\Coim{\mathop{\mathrm{Coim}}}

% Boxtimes
%
\DeclareMathSymbol{\boxtimes}{\mathbin}{AMSa}{"02}

%
% Macros for moduli stacks/spaces
%
\def\QCohstack{\mathcal{QC}\!\mathit{oh}}
\def\Cohstack{\mathcal{C}\!\mathit{oh}}
\def\Spacesstack{\mathcal{S}\!\mathit{paces}}
\def\Quotfunctor{\mathrm{Quot}}
\def\Hilbfunctor{\mathrm{Hilb}}
\def\Curvesstack{\mathcal{C}\!\mathit{urves}}
\def\Polarizedstack{\mathcal{P}\!\mathit{olarized}}
\def\Complexesstack{\mathcal{C}\!\mathit{omplexes}}
% \Pic is the operator that assigns to X its picard group, usage \Pic(X)
% \Picardstack_{X/B} denotes the Picard stack of X over B
% \Picardfunctor_{X/B} denotes the Picard functor of X over B
\def\Pic{\mathop{\mathrm{Pic}}\nolimits}
\def\Picardstack{\mathcal{P}\!\mathit{ic}}
\def\Picardfunctor{\mathrm{Pic}}
\def\Deformationcategory{\mathcal{D}\!\mathit{ef}}
\usepackage{cleveref}
\crefname{originallemma}{original lemma}{original lemmas}
\Crefname{originallemma}{Original lemma}{Original lemmas}
\makeatletter
\newlabel{lemma-stalk-exact}{{03PT}{source}{}{sourcepointer-1}{}}
\newlabel{lemma-stalk-exact@cref}{{[originallemma][0][]03PT}{[1][0][]source}{}{}{}}
\makeatother
\crefname{originallemma}{original lemma}{original lemmas}
\Crefname{originallemma}{Original lemma}{Original lemmas}
\makeatletter
\newlabel{sites-section-sheaves-injective}{{00WL}{source}{}{sourcepointer-2}{}}
\newlabel{sites-section-sheaves-injective@cref}{{[originallemma][0][]00WL}{[1][0][]source}{}{}{}}
\makeatother
\title{Exactness of maps and sequences of sheaves tested on stalks}
\author{The Stacks Project Authors}
\date{stacks-project:etale-cohomology.tex:local-snapshot-20261011; GNU FDL 1.2 or later}
\begin{document}
\maketitle
\noindent\textit{Editorial note.} All delivered mathematical text is the author's own, copied from the pinned Stacks Project chapter file etale-cohomology.tex; no statement or proof step was rewritten, reordered or removed.\par\medskip
\noindent\textit{Editorial note.} The proof is the author's complete argument for both assertions, the sequence case being deduced from the map case inside the counted span.\par\medskip
\noindent\textit{Editorial note.} Every labelled reference occurring inside the counted spans is carried in assembly.external\_references as a non-load-bearing navigation pointer (same-chapter ones at their definition site, cross-chapter ones at the invocation inside the counted span), so no load-bearing mathematics is replaced by a stub. No bibliography citation occurs inside the counted spans.\par\medskip
\noindent\textit{Editorial note.} The checked chapter list for same-name collisions is the set of stacks chapters staged in this repository: chow, cohomology, divisors, etale-cohomology, flat, more-algebra, more-morphisms, morphisms, spaces-more-morphisms, varieties together with this card's chapter etale-cohomology.\par\medskip
\noindent\textit{Editorial note.} The preamble is the chapter's own native preamble excerpt already used for this source: the Xy-pic setup (native preamble lines 14-18) and the native theorem declarations with the mathematical macros (native preamble lines 154-217).\par\medskip
\noindent\textit{Allowed standard prerequisites (editorial).}\begin{itemize}
\item Stalk of a sheaf at a geometric point. Same chapter, the stalk functor at a geometric point and its exactness properties. Supplies the functors on which exactness is tested.
\item Criterion for a map of sheaves to be injective or surjective. Same chapter, the description of injections and surjections of sheaves by their behaviour on sections over a cofinal system of coverings. Supplies the argument showing that stalkwise injectivity or surjectivity implies the corresponding property of the map.
\end{itemize}\medskip
\ifdefined\hypertarget\hypertarget{sourcepointer-1}{}\fi
\noindent\textit{Source reference.} Original lemma 03PT: sources/stacks-etale-cohomology/pinned-original-source.tex, source lines 3456--3456. Navigation pointer invoked at proof line 3558; the counted span contains the pointer only and the cited result is not reproduced inside it, so the pointer is registered as a non-load-bearing navigation stub in the established form for Stacks material. No bibliography citation is involved.\par\medskip
\ifdefined\hypertarget\hypertarget{sourcepointer-2}{}\fi
\noindent\textit{Source reference.} Original lemma 00WL: sources/stacks-etale-cohomology/pinned-original-source.tex, source lines 3579--3579. Navigation pointer invoked at proof line 3579; the counted span contains the pointer only and the cited result is not reproduced inside it, so the pointer is registered as a non-load-bearing navigation stub in the established form for Stacks material. No bibliography citation is involved.\par\medskip
\begin{theorem}
\label{theorem-exactness-stalks}
Let $S$ be a scheme.
A map $a : \mathcal{F} \to \mathcal{G}$ of sheaves of sets is injective
(resp.\ surjective) if and only if the map on stalks
$a_{\overline{s}} : \mathcal{F}_{\overline{s}} \to \mathcal{G}_{\overline{s}}$
is injective (resp.\ surjective) for all geometric points of $S$.
A sequence of abelian sheaves on $S_\etale$ is exact
if and only if it is exact on all stalks at geometric points of $S$.
\end{theorem}

\begin{proof}
The necessity of exactness on stalks follows from
Lemma \ref{lemma-stalk-exact}.
For the converse, it suffices to show that a map of sheaves is surjective
(respectively injective) if and only if it is surjective (respectively
injective) on all stalks. We prove this in the case of surjectivity, and omit
the proof in the case of injectivity.

\medskip\noindent
Let $\alpha : \mathcal{F} \to \mathcal{G}$ be a map of sheaves such
that $\mathcal{F}_{\overline{s}} \to \mathcal{G}_{\overline{s}}$
is surjective for all geometric points. Fix
$U \in \Ob(S_\etale)$
and $s \in \mathcal{G}(U)$. For every $u \in U$ choose some
$\overline{u} \to U$ lying over $u$ and an \'etale neighborhood
$(V_u , \overline{v}_u) \to (U, \overline{u})$ such that
$s|_{V_u} = \alpha(s_{V_u})$ for some
$s_{V_u} \in \mathcal{F}(V_u)$.
This is possible since $\alpha$ is surjective on
stalks. Then $\{V_u \to U\}_{u \in U}$
is an \'etale covering on which the restrictions of $s$
are in the image of the map $\alpha$.
Thus, $\alpha$ is surjective, see
Sites, Section \ref{sites-section-sheaves-injective}.
\end{proof}

\end{document}
